% !TEX program = pdflatex
%
% A cold-neutron slope discriminator between discrete-substrate and
% environmental decoherence at ILL Grenoble.
%
% PRL-style letter (revtex4-2), two-column, self-contained bibliography.
% Companion documents:
%   - Letter-ColdNeutron-Slope-Test.md (source / human-readable version)
%   - Letter-ColdNeutron-SupplementaryMethods.md (experimental details)
%   - Letter-ColdNeutron-CoverLetter.md (editor cover letter)
%
% If revtex4-2 is unavailable the preamble can be flipped to
% \documentclass[twocolumn,10pt]{article} with minimal downstream edits;
% the PRL-specific features used here are only \affiliation, twocolumn,
% and the thebibliography block.

\documentclass[aps,prl,twocolumn,superscriptaddress,nofootinbib,10pt]{revtex4-2}

\usepackage[T1]{fontenc}
\usepackage[utf8]{inputenc}
\usepackage{amsmath,amssymb,amsthm,mathtools}
\usepackage{bm}
\usepackage{booktabs}
\usepackage{hyperref}
\hypersetup{hidelinks}

% ---------- Shorthands ----------
\newcommand{\lP}{\ell_P}
\newcommand{\tP}{t_P}
\newcommand{\EP}{E_P}
\newcommand{\eps}{\varepsilon}
\newcommand{\dcomp}{\delta_{\mathrm{comp}}}
\newcommand{\Nmax}{N_{\max}}
\newcommand{\mathlibv}{\textsc{Mathlib}~v4.29}

\begin{document}

\title{A cold-neutron slope discriminator between discrete-substrate\\
  and environmental decoherence at ILL Grenoble}

\author{Norbert Marchewka}
\email{norbert.marchewka44@gmail.com}
\affiliation{Independent researcher, Poland}

\author{the OmegaTheory V2 Collaboration}
\affiliation{Machine-checked formalisation, Lean 4 / \mathlibv}

\date{2026-04-15, draft v1.0}

\begin{abstract}
We propose a tabletop slope discriminator---$-\log V \propto 1/v$
vs.\ $1/v^{2}$---that distinguishes the OmegaTheory discrete-spacetime
substrate from environmental (thermal-bath) decoherence at current
ILL Grenoble VCN sensitivity.  The substrate thesis predicts that the
log-visibility of a cold-neutron interferometer, at fixed arm length
$L$ and swept beam velocity $v$, scales \emph{linearly} in $1/v$, via
the closed-form identity $(d/v)\,c\,\eps(N)$ derived from a single
Planck-tick truncation error.  Lindblad-type thermal-bath decoherence
instead gives slope $\geq +2$.  The test is a \emph{shape
discriminator}: it is independent of the absolute noise amplitude and
robust against detector-floor calibration.  The existing ILL PF2-VCN
program, commissioned 2025 with nanodiamond--polymer gratings, already
has the velocity range ($v \in [30, 700]\,\mathrm{m/s}$, $\sim 1.37$
decades) needed to resolve slope $+1$ from slope $+2$ at
$\gtrsim 10\sigma$ with an $11.5$-beam-day campaign at
$\sim\$97.5$K.  A positive result would corroborate the first
verified OmegaTheory prediction~\cite{Huang2024Diraq} in a second,
orthogonal kinematic channel.
\end{abstract}

\maketitle

\section{Background}

OmegaTheory V2 is a discrete-spacetime substrate in which
four-dimensional geometry, quantum mechanics, and thermodynamics
emerge from a local lattice update rule at the Planck
scale~\cite{Marchewka2026Main}. The central quantitative input is a
single dimensionless \emph{irrationality truncation error}
\begin{equation}
  \eps(N) \;=\; \frac{4}{2N+3},
\end{equation}
derived from the finite continued-fraction truncation of $\pi$, $e$
and $\sqrt{2}$ in the substrate's action
density~\cite{Marchewka2026AppK}. The per-tick computational
uncertainty
\begin{equation}
  \dcomp(N) \;=\; \lP\,\eps(N)
\end{equation}
accumulates linearly in the tick count $K$ during any physical
evolution, producing an eight-item catalog of falsifiable
predictions~\cite{Marchewka2026AppJ}. At finite temperature $T$ the
admissible iteration budget is capped by
$\Nmax(T) = \hbar/(k_B\,T\,\tP)$, giving all temperature-dependent
observables an explicit closed form.

\textit{Verified prediction (first of the catalog).}---The substrate
predicts a \emph{power-law} thermal scaling of qubit coherence times
rather than Arrhenius activation.  Huang \emph{et al.}~\cite{Huang2024Diraq}
reported $T_1 \propto T^{-2.0}\ldots T^{-3.1}$ and $T_2 \propto T^{-1.0}\ldots T^{-1.1}$
for silicon spin qubits above $1\,$K, ruling out Arrhenius by
$\sim 48$ orders of magnitude. OmegaTheory's $F(T) = F_0/(1+\alpha T)$
predicts this shape.

The remaining catalog entries are either below current sensitivity by
many orders of magnitude (atomic-clock floor, UHECR dispersion,
spin-flip rate) or require platforms that do not yet exist
(mesoscopic matter-wave gravitational decoherence).
\textbf{The one exception}---and the subject of this letter---is the
cold-neutron slope test, for which the enabling hardware (ILL Grenoble
VCN~\cite{Ackermann2026VCN}) is operating \emph{now}.

\section{The prediction}

\subsection{Substrate slope}
Consider a two-path interferometer of arm length $L$ traversed by a
neutron of velocity $v$. The number of substrate ticks along each arm
is $K = \lfloor L/(v\,\tP)\rfloor$. The accumulated per-channel
infidelity is bounded above by the substrate distance--velocity
identity
\begin{equation}
  \frac{d}{v\,\tP}\,\dcomp(N) \;=\; \frac{d}{v}\,c\,\eps(N),
\end{equation}
where the factor $\lP/\tP = c$ (from $\tP \equiv \lP/c$) cancels
$\lP$ exactly. This is Lean theorem
\texttt{teleportation\_distance\_velocity\_identity} and its
floored-$K$ bound \texttt{teleportation\_distance\_velocity\_bound},
both in \texttt{OmegaTheory/Predictions/}
\texttt{StochasticTeleportation.lean}, Lean 4 / \mathlibv, no
\texttt{sorry}, no new axioms.  The log-visibility thus scales as
\begin{equation}
  -\log V_{\mathrm{sub}}(L, v) \;\propto\;
    \frac{L}{v}\,\frac{c\,\eps(N)}{\hbar}
    \quad \text{(substrate, shape)}.
\end{equation}
At fixed $L$, swept $v$, the log--log plot of $-\log V$ vs $1/v$ has
\emph{slope} $+1$.  Monotonicity in $1/v$ is
\texttt{slope\_distinguisher\_inv\_v} (same file), with the
qualitative statement: faster transit, fewer ticks, smaller substrate
infidelity contribution.  All four theorems assume only positivity of
$v, \tP, c$ (proved from \texttt{Spacetime/Constants.lean}); no axiom
beyond the four Planck-unit constants is introduced.

\subsection{Decoherence slope}
Standard thermal-bath / Lindblad decoherence~\cite{Diosi1984,Penrose1996}
couples the flight time $\tau = L/v$ quadratically or more strongly
to the visibility:
\begin{equation}
  -\log V_{\mathrm{bath}}(L, v) \;\propto\;
  \left(\frac{L}{v}\right)^{2} \Gamma_{\mathrm{bath}}(T)
  \quad \text{(thermal bath)}.
\end{equation}
The quadratic-or-higher exponent on $1/v$ arises from the lowest-order
second cumulant of a stationary noise process seen by the particle
during transit; sub-Ohmic baths and correlated noise give larger
exponents. Any $v$-scaling weaker than $1/v^{2}$ is inconsistent with
every stationary environmental noise model we are aware of.  This is
the discriminator's logical backbone: \emph{slope $+1$ cannot be
reproduced by thermal-bath decoherence, only by a per-tick
substrate}.

\subsection{Why shape and not amplitude}
The substrate's $\eps(N)$ is of order $4/(2\Nmax+3)$, which for
lab temperatures ($T\sim 10\,$mK) gives a per-tick infidelity
$\sim 10^{-52}$ [\texttt{teleportation\_fidelity\_substrate\_bound\_low\_T},
closed form $K\cdot 2\lP k_B T/\EP$]. The \emph{amplitude} is thus
too small for direct detection.  What \emph{is} measurable is the
\emph{slope of $-\log V$ against $1/v$} at fixed $L$.  The slope is
scale-invariant under overall amplitude rescaling, and the substrate's
$\lP \to c$ cancellation guarantees that the coefficient of $1/v$ is
of physical order $(L\,c\,\eps)/\hbar$, not Planck-suppressed.  A
detector that can measure any visibility above its floor can in
principle read off the slope.  What it cannot do is distinguish
substrate from bath if they are not functionally separated, and this
letter's point is that \emph{they are}, by the integer $1$ vs.\ $2$
gap in the log--log exponent.

\section{Experimental protocol}

\subsection{Platform}
\emph{ILL Grenoble PF2-VCN}---Very Cold Neutron beamline, Institut
Laue-Langevin; commissioned 2025 (May--October) with
nanodiamond--polymer composite gratings~\cite{Ackermann2026VCN}.  The
PF2-VCN configuration delivers neutrons in $v \in [30, 700]\,\mathrm{m/s}$
(equivalently $\lambda \in [5.6, 132]\,\mathrm{\AA}$), spanning $\sim 1.37$
velocity decades.  Analogous capability exists at NIST NCNR and
PSI-Villigen (ULTRA); ILL is our first-choice platform because of
the existing VCN program and the demonstrated nanodiamond gratings.

\subsection{Parameter sweep}
\begin{itemize}\itemsep0pt
  \item \textbf{Arm lengths.} $L \in [0.3, 3]\,$m (three points:
  $0.3, 1, 3\,$m; native Ackermann grating pitch covers this
  directly, Supp.~Methods \S S1.1~\cite{Alnitak2026Supp}). Longer
  arms to $\sim 10\,$m accessible via custom gratings at
  $\$10$--$20$K per arm.
  \item \textbf{Velocities.} $v \in [30, 700]\,$m/s, selected by a
  disk chopper at the interferometer entry; 10--12 values
  logarithmically spaced.
\end{itemize}

\subsection{Measurement}
The observable is visibility
$V(L,v) = (I_{\max} - I_{\min})/(I_{\max} + I_{\min})$,
measured from the first-harmonic amplitude of the Talbot--Lau fringe
pattern after a path-length phase scan. For each fixed $L$, form the
log--log plot $\log[-\log V(L,v)]$ vs $\log(1/v)$ and fit a linear
slope. Repeat for each $L$.

\paragraph*{Substrate-consistent result:} slopes all fall in
$[0.9, 1.1]$ with per-$L$ uncertainties compatible.
\paragraph*{Thermal-bath result:} slopes all fall in
$[1.7, 2.5]$.

\subsection{Sensitivity budget}
Current neutron-interferometer visibility floors are instrument- and
technique-dependent; atom interferometry currently reaches
$V \sim 10^{-9}$.  The cold-neutron VCN first-harmonic visibility
number is not published in the Ackermann \emph{et al.}\ abstract or
accessible preprint sections; a nominal operating visibility
$V_{\max} \sim 0.3$--$0.6$ is plausible for a Talbot--Lau
nanodiamond-grating configuration at VCN wavelengths
(Supp.~Methods \S S1.3~\cite{Alnitak2026Supp}) [TBD: confirm with
ILL PF2 team or Ackermann follow-up].  Slope discrimination at $+1$
vs $+2$ requires only that the \emph{dynamic range} of $-\log V$ be
resolvable above the instrument noise; since the slope is scale-free,
the absolute visibility need not be small---it need only be a
function of $v$.

Statistical slope uncertainty for a $n$-point linear fit in log--log
space goes as
\begin{equation}
  \sigma_{\mathrm{slope}} \;\approx\;
  \frac{\sigma_{\log[-\log V]}}{\sqrt{n\,\mathrm{Var}(\log(1/v))}}.
\end{equation}
For $n = 10$ points spanning $1.37$ decades
($\mathrm{Var}(\log(1/v)) \approx 0.18$) and a per-point log-visibility
error $\sigma \approx 0.1$ (typical for $\sim 10^{4}$ counts per
point), $\sigma_{\mathrm{slope}} \approx 0.07$.  The integer gap
$|\mathrm{slope}_{\mathrm{sub}} - \mathrm{slope}_{\mathrm{bath}}| = 1$
is thus $\sim 14\sigma$: statistically easy.  The full systematics
budget (neutron-gas scattering, gravitational phase-shift dephasing,
grating thermal stability) is in Supp.~Methods \S S3.

\subsection{Cost}
Per Supp.~Methods \S S4.4 Table S5~\cite{Alnitak2026Supp}: total
campaign $11.5$ beam-days at $\sim\$97.5$K, broken down as
calibration (1.5\,d, \$7.5K), velocity scan (5.5\,d, \$27.5K),
length cross-check (2.5\,d, \$12.5K), custom nDPC grating
fabrication for two off-nominal arm lengths (\$20K),
turbo-pumped vacuum-enclosure retrofit (\$15K, hard prerequisite
replacing the Ackermann commissioning He-filled box),
contingency (2\,d, \$15K).  The enabling nanodiamond--polymer
gratings already exist at PF2-VCN.

\section{Falsification criterion}

\begin{table}[htbp]
  \centering
  \renewcommand{\arraystretch}{1.1}
  \begin{tabular}{@{}lp{0.6\linewidth}@{}}
    \toprule
    Slope at 95\% CL & Interpretation \\
    \midrule
    $[0.9, 1.1]$ & Substrate confirmed. Slope $+1$ is inaccessible
    to every stationary thermal-bath model we know. \\
    $\geq 1.7$ & Substrate falsified in the cold-neutron channel;
    thermal-bath decoherence favoured. \\
    $[1.1, 1.7]$ & Inconclusive; requires wider $v$ range ($v$ down
    to $\lesssim 30\,$m/s via PF2-UCN) or longer integration. \\
    \bottomrule
  \end{tabular}
  \caption{Fit outcome classification.}
  \label{tab:outcome}
\end{table}

\textit{Shape survives amplitude uncertainty.}---Slope fitting is
insensitive to any single-prefactor calibration error, because
$V \to \kappa V$ adds a constant to $\log[-\log V]$ without changing
the fitted slope.  It \emph{is} sensitive to multiplicative
$v$-dependent systematics---e.g.\ a detector efficiency that scales
with neutron velocity.  Standard neutron-interferometry technique
(reversible coherent beam rotation, incoherent beam calibration)
makes this tractable.

\section{Theoretical significance}

The cold-neutron slope test is the first cleanly functional (not
amplitude-limited) discriminator between the OmegaTheory substrate
and environmental decoherence. Its significance is best understood in
relation to the already-verified Diraq result:
\begin{itemize}\itemsep0pt
  \item \textbf{Diraq 2024~\cite{Huang2024Diraq}:} verified the
  $T$-scaling \emph{shape} (power-law, not Arrhenius); ruled out
  Arrhenius by $\sim 48$ OOM.
  \item \textbf{Cold-neutron slope (this proposal):} tests the
  $v$-scaling \emph{shape} (linear, not quadratic-or-higher in
  $1/v$); would rule out thermal-bath decoherence by an integer
  shape exponent.
\end{itemize}
If both pass, two independent substrate shape predictions---operating
on physically orthogonal kinematic variables ($T$, the thermal budget;
$v$, the transit-time budget)---are verified from two different
physical systems (silicon spin qubits, cold neutrons). The substrate's
underlying mechanism is the per-tick truncation error $\eps(N)$, and
the same $\eps(N)$ governs both $T$-scaling (via
$\Nmax(T) \propto 1/T$) and $v$-scaling (via
$K = d/(v\tP) \propto 1/v$).

Equally important is the \emph{asymmetry of falsification}: the slope
test is cheap and integer-gap decisive.  A single negative result
(slope $\geq 1.7$) rules out this branch of the substrate hypothesis
in the cold-neutron channel, and the theory would have to survive on
the $T$-scaling verification alone.

\section{Conclusion}

We have proposed a cold-neutron interferometric slope test at ILL
Grenoble VCN as the first feasible experimental discriminator between
the OmegaTheory discrete substrate and environmental decoherence. The
prediction is a machine-checked theorem
(\texttt{OmegaTheory.Predictions.teleportation\_distance\_velocity\_identity},
Lean 4 / \mathlibv; 0 \texttt{sorry}, 0 new axioms beyond the
Planck-unit constants) that the log-visibility of a two-path
interferometer at fixed arm length scales as $1/v$ (substrate) rather
than $1/v^{2}$ or steeper (thermal bath). The test is robust against
absolute-amplitude calibration, feasible at $\sim\$97.5$K of marginal
beam time~\cite{Alnitak2026Supp}, and complementary to the
already-verified Diraq anti-Arrhenius signature~\cite{Huang2024Diraq}.
A positive outcome would constitute a second, orthogonal corroboration
of the substrate's central prediction mechanism; a negative outcome
would place a clean, shape-level constraint on the discrete-spacetime
hypothesis.

% ----------------------------------------------------------------------
% Lean-source footer: the primary formal evidence underpinning the
% prediction.  These files are the authoritative source; citations above
% are the human-readable layer over them.
% ----------------------------------------------------------------------

\paragraph*{Lean source files (primary formal evidence).}
Public repository permalink: [TBD: permalink to be added on submission].
\begin{itemize}\itemsep0pt\small
  \item \texttt{OmegaTheory/Predictions/StochasticTeleportation.lean}
  (Regulus, 2026-04-15):
  \texttt{teleportation\_distance\_velocity\_identity},
  \texttt{slope\_distinguisher\_inv\_v},
  \texttt{teleportation\_fidelity\_substrate\_bound\_low\_T}.
  \item \texttt{OmegaTheory/Emergence/MassAsDelay.lean} (Aldebaran,
  2026-04-15): \texttt{perTickDelay\_high\_momentum\_bound}.
  \item \texttt{OmegaTheory/Emergence/SnapshotPropagator.lean}:
  \texttt{accumulatedSnapshotError\_add} (linear cumulation).
  \item \texttt{OmegaTheory/Irrationality/Uncertainty.lean}:
  \texttt{iterationBudget\_decreases\_with\_T},
  \texttt{computationalUncertainty\_decreasing}.
\end{itemize}
All files built clean under Lean~4 / \mathlibv\ at the
\textbf{3{,}835-jobs-GREEN, 0-\texttt{sorry}, 8-physical-axiom,
8{,}996-\texttt{Theorem}-node top-level state} (2026-04-21).

\textit{Draft prepared by agent Bellatrix ($\gamma$ Orionis),
OmegaTheory V2 paper-submission track, wave 4 $\cdot$ 2026-04-15.}

% ----------------------------------------------------------------------
% Bibliography (self-contained \bibitem blocks; no external .bib needed).
% ----------------------------------------------------------------------

\begin{thebibliography}{9}

\bibitem{Huang2024Diraq}
J.~Y.~Huang, R.~Y.~Su, W.~H.~Lim \emph{et al.},
``High-fidelity spin qubit operation and algorithmic initialization
above 1~K,'' \emph{Nature} \textbf{627}, 772--777 (2024).
doi:\href{https://doi.org/10.1038/s41586-024-07160-2}{10.1038/s41586-024-07160-2}.
Verified Arrhenius rule-out; first OmegaTheory prediction confirmed
in published experiment.

\bibitem{Ackermann2026VCN}
M.~Ackermann \emph{et al.},
``Very-cold-neutron interferometry with nanodiamond-polymer composite
gratings at ILL-PF2-VCN,'' arXiv:2604.09312 (2026).
ILL VCN commissioning report; platform for the proposed test.

\bibitem{PierreAuger2026LIV}
Pierre Auger Collaboration,
``Constraints on hadronic Lorentz invariance violation from muon
fluctuations in cosmic-ray air showers,'' arXiv:2602.14720 (Feb.\ 2026).
Current strongest hadronic-sector LIV bounds; substrate predictions
consistent with all present null results.

\bibitem{Diosi1984}
L.~Di\'osi,
``Gravitation and quantum-mechanical localization of macro-objects,''
\emph{Phys.\ Lett.\ A} \textbf{105}, 199--202 (1984).
doi:\href{https://doi.org/10.1016/0375-9601(84)90397-9}{10.1016/0375-9601(84)90397-9}.

\bibitem{Penrose1996}
R.~Penrose,
``On gravity's role in quantum state reduction,''
\emph{Gen.\ Rel.\ Grav.} \textbf{28}, 581--600 (1996).
doi:\href{https://doi.org/10.1007/BF02105068}{10.1007/BF02105068}.
Comparison models for environmental gravitational decoherence.

\bibitem{Marchewka2026Main}
N.~Marchewka,
``OmegaTheory V2: emergent spacetime, quantum mechanics, and
thermodynamics from an irrationality-truncated lattice substrate,''
paper draft (2026).
Substrate derivation, grand QM emergence theorem, full Lean~4
formalisation (\mathlibv; 3{,}835 build jobs GREEN, 0~\texttt{sorry},
8 physical axioms; 8{,}996 \texttt{:Theorem} nodes in the
\texttt{OmegaTheoryV2} Neo4j graph as of 2026-04-21).

\bibitem{Marchewka2026AppJ}
N.~Marchewka,
``Appendix J: Comprehensive Experimental Catalog (Consolidated),''
OmegaTheory V2 supplementary material (2026-04-15).
Full prediction catalog; \S1.1 is the extended version of the test
proposed in this letter.

\bibitem{Marchewka2026AppK}
N.~Marchewka,
``Appendix K: Irrationality Genesis of Predictions,''
OmegaTheory V2 supplementary material (2026).
Derivation of $\eps(N) = 4/(2N+3)$ from continued-fraction truncation
of $\pi, e, \sqrt{2}$.

\bibitem{Alnitak2026Supp}
Alnitak (OmegaTheory V2 Collaboration),
``Supplementary Methods: Experimental sensitivity and noise budget
for the cold-neutron slope test at ILL Grenoble PF2-VCN,''
companion to this letter (2026-04-15).
File: \texttt{Letter-ColdNeutron-SupplementaryMethods.md}.
Full signal/systematics/beam-time budget; $11.5$\,d, \$97.5K campaign.

\end{thebibliography}

\end{document}
